\begin{table}[t]
\centering
\small
\begin{tabular}{llcccc}
\toprule
strategy & norm & synthetic & real & real & vacuous \\
         &      &           & (unweighted) & (weighted) & (\%) \\
\midrule
naive ($\bar\theta$) & raw & 0.780 & 0.971 & 0.971 & 0.0  \quad$\leftarrow$ baseline \\
 & $v_y$ & 0.743 & 0.941 & 0.941 & 0.0 \\
 & $v_y + v_c$ & 0.765 & 0.971 & 0.912 & 0.0 \\
 & $\delta_c,\delta_y$ & 0.770 & 0.971 & 0.882 & 0.0 \\
\midrule
naive + $\Delta_c$ & raw & 0.948 & 1.000 & 1.000 & 0.0 \\
 & $v_y$ & 0.950 & 1.000 & 1.000 & 0.0 \\
 & $v_y + v_c$ & 0.963 & 1.000 & 1.000 & 0.0 \\
 & $\delta_c,\delta_y$ & 0.980 & 1.000 & 1.000 & 0.0 \\
\midrule
surrogate ($\theta^*$) & raw & 0.765 & 0.971 & 0.824 & 0.0 \\
 & $v_y$ & 0.725 & 0.941 & 0.912 & 0.0 \\
 & $v_y + v_c$ & 0.728 & 0.941 & 0.912 & 0.0 \\
 & $\delta_c,\delta_y$ & 0.738 & 0.941 & 0.882 & 0.0 \\
\bottomrule
\end{tabular}
\caption{Joint coverage at nominal $1-\alpha = 0.68$ (Bonferroni-corrected across the $d=4$ coordinates). Synthetic test $n=400$, real test $n=34$. The first row is standard split conformal prediction calibrated on the data-generating parameters $\bar\theta$, without surrogate labels or covariate-shift reweighting. ``vacuous'' is the fraction of real test systems whose interval is unbounded.}
\label{tab:cp-ablation}
\end{table}
